\documentclass[11pt,a4paper]{article}

\usepackage[margin=2.15cm]{geometry}
\usepackage[turkish]{babel}
\usepackage{fontspec}
\setmainfont{lmroman10-regular.otf}[BoldFont=lmroman10-bold.otf]
\usepackage{microtype}
\usepackage{booktabs}
\usepackage{longtable}
\usepackage{xcolor}
\usepackage{listings}
\usepackage{fancyhdr}
\usepackage{hyperref}
\usepackage{enumitem}

\hypersetup{
  colorlinks=true,
  linkcolor=blue!45!black,
  urlcolor=blue!55!black,
  pdftitle={NUCLEO-G491RE BLDC Komutasyon Simülasyonu},
  pdfauthor={Çalışma Kılavuzu}
}
\definecolor{stmblue}{RGB}{0,103,177}
\definecolor{codebg}{RGB}{245,247,250}
\lstset{
  basicstyle=\ttfamily\footnotesize,
  backgroundcolor=\color{codebg},
  frame=single,
  rulecolor=\color{black!20},
  breaklines=true,
  columns=fullflexible,
  showstringspaces=false
}
\pagestyle{fancy}
\fancyhf{}
\lhead{\textcolor{stmblue}{NUCLEO-G491RE | BLDC UART Simülasyonu}}
\rhead{Çalışma Kılavuzu}
\cfoot{\thepage}
\setlength{\headheight}{14pt}
\setlist[itemize]{leftmargin=*,noitemsep,topsep=3pt}
\setlist[enumerate]{leftmargin=*,topsep=3pt}

\begin{document}
\pagecolor{white}
\color{black}

\begin{titlepage}
  \centering
  \vspace*{2.4cm}
  {\Huge\bfseries\textcolor{stmblue}{NUCLEO-G491RE}\par}
  \vspace{0.35cm}
  {\LARGE 3 Fazlı BLDC Komutasyon Simülasyonu\par}
  \vspace{0.25cm}
  {\Large STM32CubeMX + STM32CubeIDE + ST-LINK VCP + CoolTerm\par}
  \vfill
  \begin{tabular}{rl}
    \textbf{Proje adı:} & \texttt{staj\_motor\_simule\_g491} \\
    \textbf{Kart:} & NUCLEO-G491RE \\
    \textbf{MCU:} & STM32G491RET6 \\
    \textbf{Seri arayüz:} & LPUART1 / ST-LINK Virtual COM Port \\
    \textbf{Tarih:} & 02 Eylül 2026 \\
  \end{tabular}
  \vfill
  \fbox{\parbox{0.83\textwidth}{\centering
  Bu doküman, projeyi sıfırdan yeniden oluşturmak, kodun mantığını anlamak,
  karta yazmak ve CoolTerm'de sonuçları yorumlamak için hazırlanmıştır.}}
  \vspace*{1.6cm}
\end{titlepage}

\tableofcontents
\newpage

\section{Amaç ve güvenlik sınırı}

Bu uygulama gerçek bir BLDC motoru döndürmez. Bunun yerine üç GPIO pini üzerinde
BLDC'nin altı adımlı (six-step) komutasyon sırasını canlandırır. Her \textbf{200 ms}
sonunda A, B ve C fazının durumu değiştirilir; aynı bilgi LPUART1 ile bilgisayara
gönderilir ve CoolTerm'de metin veya grafik olarak izlenir.

\begin{center}
\fcolorbox{red!55!black}{red!5}{\parbox{0.88\textwidth}{
\textbf{Güvenlik:} PA4, PB4 ve PB5 yalnızca 3.3 V mantık çıkışlarıdır. BLDC motoru,
motor fazını veya MOSFET kapısını bu pinlere doğrudan bağlama. Gerçek sistemde
3-faz inverter/gate driver, akım ölçümü, koruma ve PWM gerekir. Bu çalışma
yalnızca komutasyon mantığını güvenli biçimde öğretir.
}}
\end{center}

Öğrenilen veri zinciri şöyledir:

\begin{lstlisting}
Komutasyon tablosu -> GPIO A/B/C -> LPUART1 TX
                   -> ST-LINK Virtual COM Port -> CoolTerm
\end{lstlisting}

\section{Gerekli donanım ve yazılım}

\subsection{Gereksinimler}
\begin{itemize}
  \item NUCLEO-G491RE geliştirme kartı
  \item Veri aktarımını destekleyen micro-USB kablo
  \item macOS çalıştıran bilgisayar
  \item STM32CubeMX, STM32CubeIDE, CoolTerm
  \item İsteğe bağlı: STM32CubeProgrammer
\end{itemize}

\subsection{Nucleo kartı neden işimizi kolaylaştırır?}

Kart üzerindeki ST-LINK V3E iki ayrı iş yapar:
\begin{enumerate}
  \item \textbf{SWD}: Firmware yazma ve hata ayıklama,
  \item \textbf{Virtual COM Port (VCP)}: UART verisini USB üzerinden Mac'e taşıma.
\end{enumerate}

Bu yüzden ek USB--UART dönüştürücü ve DFU modu gerekmez.

\begin{lstlisting}
Mac + CoolTerm
      |  micro-USB veri kablosu
      v
NUCLEO üzerindeki ST-LINK V3E
      |-- SWD <-> PA13 (SWDIO), PA14 (SWCLK): programlama/debug
      |-- VCP <-> PA2 (LPUART1_TX), PA3 (LPUART1_RX): UART
      v
STM32G491RET6
\end{lstlisting}

\section{STM32CubeMX ile sıfırdan proje oluşturma}

\subsection{Kart seçimi}
\begin{enumerate}
  \item STM32CubeMX'i aç.
  \item \textit{Access to Board Selector} seçeneğini seç.
  \item Arama alanına \texttt{NUCLEO-G491RE} yaz.
  \item Kartı seç ve \textit{Start Project} düğmesine bas.
\end{enumerate}

Kart seçicisini kullanmak önemlidir: CubeMX, kartın ST-LINK ve Virtual COM Port
bağlantılarını tanır. Yalnız MCU seçilirse fiziksel kart bağlantılarını senin
doğrulaman gerekir.

\subsection{LPUART alanı kilitli veya gri görünürse}

Kart şablonu bazı pinleri önceden ayırmış olabilir. Sıfırdan açılmış projede
LPUART1 modu seçilemiyorsa \textit{Pinout} menüsünden \textit{Clear Pinout} veya
ilgili \textit{Reset Configuration} işlemini kullanıp aşağıdaki ayarları yeniden
oluştur. Bu işlem mevcut projede rastgele yapılmamalıdır; önceki pin ayarlarını
silebilir.

\subsection{SYS: Serial Wire debug}

\textit{System Core > SYS} ekranında:

\begin{center}
\fbox{\texttt{Debug = Serial Wire}}
\end{center}

Bu seçim SWDIO (PA13) ve SWCLK (PA14) pinlerini ST-LINK programlayıcısına ayırır.
Bu pinleri başka GPIO işlevi için kullanma. \texttt{No Debug} seçimi, debug ve
tekrar programlamayı gereksiz derecede zorlaştırır.

\subsection{LPUART1: bilgisayara seri veri gönderme}

\textit{Connectivity > LPUART1} ekranında aşağıdaki ayarları yap:

\begin{center}
\begin{tabular}{ll}
\toprule
Alan & Değer \\
\midrule
Mode & Asynchronous \\
Baud Rate & 115200 \\
Word Length & 8 Bits \\
Parity & None \\
Stop Bits & 1 \\
Hardware Flow Control & Disable / None \\
Communication Mode & TX/RX \\
\bottomrule
\end{tabular}
\end{center}

Pinout ekranında şu eşleme görünmelidir:

\begin{center}
\begin{tabular}{lll}
\toprule
Pin & Fonksiyon & Anlamı \\
\midrule
PA2 & \texttt{LPUART1\_TX} & STM32 $\rightarrow$ ST-LINK VCP $\rightarrow$ Mac \\
PA3 & \texttt{LPUART1\_RX} & Mac $\rightarrow$ ST-LINK VCP $\rightarrow$ STM32 \\
\bottomrule
\end{tabular}
\end{center}

\textbf{TX} (Transmit), STM32'nin dışarı gönderdiği hattır; bu projede STEP
satırları PA2 üzerinden gider. \textbf{RX} (Receive), STM32'nin dışarıdan aldığı
hattır. Henüz komut almıyoruz fakat PA3'ü ayarlamak, ileride bilgisayardan
başlat/durdur/hız komutu göndermek için altyapı hazırlar.

\textbf{Neden 115200 8-N-1?} 115200, terminal uygulamalarında yaygın ve bu proje
için yeterince hızlıdır. 8-N-1; 8 veri biti, parity yok, 1 stop biti demektir.
CoolTerm tarafında da aynı seçilmelidir. CTS/RTS kullanmadığımız için Flow Control
kapalıdır.

\subsection{Üç faz GPIO çıkışı}

Pinout görünümünde şu üç pini \texttt{GPIO\_Output} yap ve etiketle:

\begin{center}
\begin{tabular}{lll}
\toprule
Pin & User Label & GPIO ayarı \\
\midrule
PA4 & \texttt{PHASE\_A} & Output Push Pull, Low, No pull, Low speed \\
PB4 & \texttt{PHASE\_B} & Output Push Pull, Low, No pull, Low speed \\
PB5 & \texttt{PHASE\_C} & Output Push Pull, Low, No pull, Low speed \\
\bottomrule
\end{tabular}
\end{center}

Bu pinler bu projede boş olan üç bağımsız GPIO'dur. \texttt{Output Push Pull},
pin HIGH iken aktif 3.3 V; LOW iken aktif GND sürmesini sağlar. Başlangıç
seviyesinin Low olması reset sonrası güvenli durum verir. \texttt{Low speed}
bu yavaş simülasyon için yeterlidir. PA5 ise karttaki LD2 kullanıcı LED'idir;
bu projede faz pini değildir.

\subsection{Clock ve kod üretimi}

CubeMX'in ürettiği yapı HSI ve PLL ile 170 MHz sistem saatini kullanır:

\begin{lstlisting}
HSI = 16 MHz
PLL girişi = 16 / 4 = 4 MHz
VCO = 4 x 85 = 340 MHz
SYSCLK = 340 / 2 = 170 MHz
\end{lstlisting}

Bu çalışma için Clock Configuration ekranında elle değişiklik yapmak gerekmez.
Project Manager ekranında şu değerleri kullan:

\begin{center}
\begin{tabular}{ll}
\toprule
Alan & Değer \\
\midrule
Project Name & \texttt{staj\_motor\_simule\_g491} \\
Toolchain / IDE & STM32CubeIDE \\
Firmware Package & STM32Cube FW\_G4 \\
Keep User Code & Etkin \\
\bottomrule
\end{tabular}
\end{center}

\textit{Keep User Code} açık olmalıdır. Bu sayede \texttt{.ioc} değişip tekrar
\textit{Generate Code} yapıldığında \texttt{USER CODE BEGIN/END} bloklarına
eklenen uygulama kodu korunur.

\section{STM32CubeIDE: kodun çalışma mantığı}

Kod \texttt{Core/Src/main.c} içindedir. Eklediğin kodu yalnızca
\texttt{USER CODE} bloklarına koy. CubeMX'in oluşturduğu
\texttt{MX\_GPIO\_Init()} ve \texttt{MX\_LPUART1\_UART\_Init()} fonksiyonlarını
silme.

\subsection{Faz durumları}

\begin{lstlisting}
typedef uint8_t PhaseState_t;

enum
{
  PHASE_FLOAT = 0U,
  PHASE_LOW   = 1U,
  PHASE_HIGH  = 2U
};
\end{lstlisting}

\begin{itemize}
  \item \texttt{PHASE\_FLOAT}: Pin sürülmez, GPIO input/yüksek empedans durumundadır.
  \item \texttt{PHASE\_LOW}: Pin aktif olarak 0 V, yani lojik 0 sürer.
  \item \texttt{PHASE\_HIGH}: Pin aktif olarak 3.3 V, yani lojik 1 sürer.
\end{itemize}

FLOAT ile LOW aynı şey değildir. FLOAT, gerçek motor sürücüsünde ilgili faz
kolunun kapatılmasına karşılık gelen lojik durumun güvenli simülasyonudur.

\subsection{Altı adımlı komutasyon tablosu}

\begin{center}
\begin{tabular}{cccc}
\toprule
Adım & PHASE\_A & PHASE\_B & PHASE\_C \\
\midrule
1 & HIGH (+1) & LOW (-1) & FLOAT (0) \\
2 & HIGH (+1) & FLOAT (0) & LOW (-1) \\
3 & FLOAT (0) & HIGH (+1) & LOW (-1) \\
4 & LOW (-1) & HIGH (+1) & FLOAT (0) \\
5 & LOW (-1) & FLOAT (0) & HIGH (+1) \\
6 & FLOAT (0) & LOW (-1) & HIGH (+1) \\
\bottomrule
\end{tabular}
\end{center}

Her adımda bir faz HIGH, bir faz LOW, diğer faz FLOAT olur. Altıncı adımdan
sonra tekrar ilk adıma dönülür.

\begin{lstlisting}
#define MOTOR_STEP_COUNT      6U
#define MOTOR_STEP_PERIOD_MS 200U

static const uint8_t commutation_table[MOTOR_STEP_COUNT] =
{
  PACK_PHASES(PHASE_HIGH,  PHASE_LOW,   PHASE_FLOAT),
  PACK_PHASES(PHASE_HIGH,  PHASE_FLOAT, PHASE_LOW),
  PACK_PHASES(PHASE_FLOAT, PHASE_HIGH,  PHASE_LOW),
  PACK_PHASES(PHASE_LOW,   PHASE_HIGH,  PHASE_FLOAT),
  PACK_PHASES(PHASE_LOW,   PHASE_FLOAT, PHASE_HIGH),
  PACK_PHASES(PHASE_FLOAT, PHASE_LOW,   PHASE_HIGH)
};
\end{lstlisting}

\texttt{PACK\_PHASES}, üç faz durumunu ikişer bit halinde tek byte'ta saklar.
\texttt{UNPACK\_PHASE} ise seçilen adımdan A/B/C bilgisini tekrar ayırır.

\subsection{Set\_Phase: HIGH, LOW ve FLOAT uygulanması}

\begin{lstlisting}
static void Set_Phase(GPIO_TypeDef *port, uint16_t pin,
                      PhaseState_t state)
{
  GPIO_InitTypeDef gpio = {0};
  gpio.Pin = pin;
  gpio.Pull = GPIO_NOPULL;
  gpio.Speed = GPIO_SPEED_FREQ_LOW;

  if (state == PHASE_FLOAT)
  {
    gpio.Mode = GPIO_MODE_INPUT;
    HAL_GPIO_Init(port, &gpio);
  }
  else
  {
    gpio.Mode = GPIO_MODE_OUTPUT_PP;
    HAL_GPIO_Init(port, &gpio);
    HAL_GPIO_WritePin(port, pin,
      (state == PHASE_HIGH) ? GPIO_PIN_SET : GPIO_PIN_RESET);
  }
}
\end{lstlisting}

\begin{itemize}
  \item HIGH: pin Output Push Pull olur, sonra \texttt{GPIO\_PIN\_SET} yazılır.
  \item LOW: pin Output Push Pull olur, sonra \texttt{GPIO\_PIN\_RESET} yazılır.
  \item FLOAT: pin input moda alınır; pin artık aktif sürülmez.
\end{itemize}

\subsection{Motor\_Advance, SysTick ve UART}

\texttt{Motor\_Advance()} güncel adımı seçer, A/B/C pinlerini günceller ve aynı
bilgiyi LPUART1 üzerinden CoolTerm'e yollar:

\begin{lstlisting}
STEP 1 | A: 1 | B: -1 | C: 0
STEP 2 | A: 1 | B: 0 | C: -1
STEP 3 | A: 0 | B: 1 | C: -1
STEP 4 | A: -1 | B: 1 | C: 0
STEP 5 | A: -1 | B: 0 | C: 1
STEP 6 | A: 0 | B: -1 | C: 1
\end{lstlisting}

\begin{lstlisting}
const uint32_t now_ms = HAL_GetTick();
if ((uint32_t)(now_ms - last_step_ms) >= MOTOR_STEP_PERIOD_MS)
{
  last_step_ms += MOTOR_STEP_PERIOD_MS;
  Motor_Advance();
}
\end{lstlisting}

\texttt{HAL\_GetTick()} milisaniye sayacıdır. Periyot 200 ms olduğu için tam
altı adımlı döngü $6 \times 200 = 1200$ ms, yani 1.2 s sürer.

\section{Derleme, karta yazma ve debug}

CubeIDE'de proje adına tıkla ve çekiç simgesine bas. Beklenen çıktı:

\begin{lstlisting}
Build Finished. 0 errors, 0 warnings.
\end{lstlisting}

Derleme sonunda şu dosya oluşur:

\begin{lstlisting}
Debug/staj_motor_simule_g491.elf
\end{lstlisting}

Kart USB ile bağlıyken CubeIDE içinden Run/Debug ile yükleme yapılabilir.
ST-LINK firmware'i SWD üzerinden yazar ve MCU'yu resetler. Kartı DFU moduna
alma gerekmez.

Eğer CubeIDE \textit{ST-Link Server is required to launch the debug session}
uyarısını verirse bu uygulama kodu hatası değildir; macOS'ta ST-LINK Server
bileşeni eksiktir. ST-LINK Server'i kurup IDE'yi yeniden aç veya
STM32CubeProgrammer'da ST-LINK/SWD bağlantısını seçip \texttt{.elf} dosyasını
programla.

\section{CoolTerm ile UART verisini görme}

CoolTerm'de \textit{Options} ekranını aç ve şu değerleri gir:

\begin{center}
\begin{tabular}{ll}
\toprule
Alan & Değer \\
\midrule
Serial Port & \texttt{/dev/cu.usbmodem...} \\
Baudrate & 115200 \\
Data Bits & 8 \\
Parity & None \\
Stop Bits & 1 \\
Flow Control & None \\
\bottomrule
\end{tabular}
\end{center}

Sondaki port numarası her bağlanışta değişebilir; önemli olan NUCLEO/ST-LINK'e
ait \texttt{/dev/cu.\allowbreak usbmodem...} girdisidir. \textit{Connect} düğmesine basınca
başlık ve STEP 1--6 verileri yaklaşık her 200 ms'de bir görünmelidir.

Metin okunmuyorsa ilk kontrol baud değeridir: hem CubeMX hem CoolTerm
\textbf{115200} olmalıdır. Sonra 8-N-1 ve Flow Control=None değerlerini kontrol
et. Port görünmüyorsa veri taşıyan USB kablosu kullan ve başka terminalin aynı
portu açık tutmadığından emin ol.

\section{CoolTerm Chart: ayar ve yorum}

\subsection{Neden dört trace oluşur?}

Okunaklı satırın içinde dört sayı vardır:

\begin{lstlisting}
STEP 1 | A: 1 | B: -1 | C: 0
     ^      ^      ^       ^
     |      |      |       +-- Trace 4: C
     |      |      +---------- Trace 3: B
     |      +----------------- Trace 2: A
     +------------------------ Trace 1: STEP numarası
\end{lstlisting}

CoolTerm satırdaki sayıları ayıklar. Bu nedenle 0--6 arasında yükselen ilk eğri
faz gerilimi değil, adım numarasıdır.

\subsection{Temiz grafik için ayar}

\begin{enumerate}
  \item \textit{View > Chart} görünümüne geç.
  \item Grafik alanında sağ tıkla, \textit{Chart Properties} seç.
  \item \textit{Only plot lines with exactly the specified number of data points}
        seçeneğini aç ve değeri \textbf{4} yap.
  \item Sağ tıkla, \textit{Trace Properties} seç.
  \item Trace 1 için \textit{Visible} seçeneğini kapat.
  \item Trace 2: \textit{Phase A}; Trace 3: \textit{Phase B}; Trace 4:
        \textit{Phase C} olarak isimlendir.
  \item \textit{Clear Data} ile eski veriyi sil; 1--2 saniye yeni veri biriktir.
\end{enumerate}

Yakınlaştırma için \textit{View > Zoom Chart > Zoom In} yolunu kullan.
Grafik üzerindeyken \texttt{Control + trackpad/mouse scroll} da zoom yapar.

\subsection{Grafiği yorumlama}

\begin{center}
\begin{tabular}{cl}
\toprule
Y değeri & Anlamı \\
\midrule
+1 & Faz GPIO'su aktif olarak 3.3 V (HIGH) sürüyor. \\
0 & Faz FLOAT; pin input/yüksek empedans, aktif sürme yok. \\
-1 & Faz GPIO'su aktif olarak 0 V/GND (LOW) sürüyor. \\
\bottomrule
\end{tabular}
\end{center}

Beklenen tekrar eden desen:

\begin{lstlisting}
A:  +1, +1,  0, -1, -1,  0, ...
B:  -1,  0, +1, +1,  0, -1, ...
C:   0, -1, -1,  0, +1, +1, ...
\end{lstlisting}

Bu eğriler sinüs değil, basamaklı/trapezoidal komutasyon sinyalleridir.
Her adımda bir faz HIGH, bir faz LOW, bir faz FLOAT görülmesi doğru sonuçtur.

\subsection{Hem okunaklı metin hem kusursuz chart istenirse}

İleride iki satır türü göndermek en temiz yöntemdir:

\begin{lstlisting}
STEP 1 | A: 1 | B: -1 | C: 0    // Terminal satırı
1,-1,0                           // Sadece chart satırı
\end{lstlisting}

Bu durumda Chart Properties'te yalnızca \textbf{3} sayı içeren satırları çiz
filtresi seçilir; grafik sadece A/B/C eğrilerinden oluşur.

\section{Kontrol listesi ve sonraki adımlar}

\subsection{Başarı kontrol listesi}
\begin{itemize}
  \item CubeIDE: \texttt{0 errors, 0 warnings}
  \item ST-LINK: Firmware karta yazılabiliyor.
  \item CoolTerm: \texttt{/dev/cu.usbmodem...}, 115200 8-N-1, Flow Control=None
  \item Terminal: STEP 1'den STEP 6'ya kadar sıra sürekli tekrar ediyor.
  \item Chart: A/B/C yalnızca -1, 0, +1 seviyelerinde.
\end{itemize}

\subsection{Geliştirme sırası}
\begin{enumerate}
  \item PA5/LD2 LED'ini her komutasyon döngüsünde değiştir.
  \item PC13/B1 butonuyla başlat/durdur ekle.
  \item UART RX ile \texttt{s}, \texttt{p}, \texttt{+}, \texttt{-} komutlarını al.
  \item \texttt{HAL\_GetTick()} yerine timer interrupt kullan.
  \item GPIO yerine timer PWM kanallarını öğren.
  \item Hall sensörü/back-EMF geri beslemesi ve gerçek gate driver ekle.
\end{enumerate}

\begin{center}
\fbox{\parbox{0.85\textwidth}{\centering
\textbf{Özet:} CubeMX donanımı tanımlar; CubeIDE HAL kodunu derler; ST-LINK
firmware'i SWD ile karta yazar; LPUART1 veriyi ST-LINK VCP üzerinden Mac'e taşır;
CoolTerm ise veriyi terminal ve grafik olarak görünür yapar.
}}
\end{center}

\end{document}
